\section{ALTERNATIVAS PROPUESTAS}
\justifying

\subsection{Antecedentes}

Existen diversos trabajos relacionados a la mejora de la enseñanza de electrónica de potencia mediante el uso de dispositivos didácticos. Por ejemplo, el trabajo ``Módulos Didácticos de Electrónica de Potencia para el control de máquinas de corriente directa'', \parencite{4915257}, desarrolla módulos que permiten estudiar características y técnicas de control de motores eléctricos. 

De manera complementaria, el ``Desarrollo de un inversor monofásico didáctico'', \parencite{SANDOVAL-INVERSOR}, se orienta al diseño de un dispositivo enfocado en una sola rama de la electrónica de potencia: la conversión de corriente continua en corriente alterna. En la misma línea, se encuentra el trabajo ``Desarrollo de una Plataforma Didáctica para Convertidores de Electrónica de Potencia Flexibles'', \parencite{VictorGuimarFra2022}, propone un sistema que facilite la comprensión y práctica de la conversión de energía a través de convertidores de distintas topologías, como el convertidor de puente completo.

Finalmente, otro trabajo que presenta una alternativa didáctica, ``Diseño e Implementación de un Prototipo de Tablero Didáctico Modular para el Laboratorio de Electrónica 3'', \parencite{ALVAREZ-TESIS}, tiene como objetivo el diseño de un tablero didáctico que facilite la labor del docente y favorezca la comprensión de los estudiantes, con un enfoque particular a la programación. Este trabajo es el único de los mencionados que fue entregado como proyecto de graduación.

\subsection{Contribución del proyecto}

Los trabajos mencionados ponen de manifiesto la necesidad de contar con un dispositivo didáctico que aproxime los contenidos teóricos a la práctica, favoreciendo el trabajo docente, la comprensión de los alumnos y con capacidad de expansión a distintos contextos académicos. Sin embargo, muchos de estos desarrollos únicamente se centran en el desarrollo de aplicaciones específicas sin abarcar otras temáticas de interés. 

Además, cabe destacar que en en la región no se dispone de ninguna solución o dispositivo similar. En este contexto, el presente proyecto busca aportar una alternativa integral, modular y accesible, que facilite la enseñanza, y que también se mantenga vigente frente a la evolución tecnológica. 


\begin{comment}
Escriba aquí con extrema claridad las diversas alternativas de solución identificadas para la implementación del proyecto y de la problemática existente.
\end{comment}